\chapter{Percurso Metodológico}

\section{Natureza, abordagem e estratégia da pesquisa}

A pesquisa será de natureza aplicada, com abordagem predominantemente qualitativa e apoio de dados quantitativos descritivos. Não buscará medir o efeito isolado da plataforma nem estabelecer relações causais generalizáveis entre a intervenção e o desempenho de estudantes. A análise examinará indícios presentes nas produções individuais e nas situações observadas durante a aplicação da sequência didática.

A estratégia principal será uma intervenção pedagógica. \citeonline{damiani2013discutindo} caracterizam esse tipo de investigação pelo planejamento e pela implementação de mudanças ou inovações pedagógicas, acompanhados da análise de seus efeitos. Essa formulação é compatível com o projeto porque os recursos da intervenção serão planejados para uma situação de ensino, aplicados no campo e examinados a partir das evidências produzidas.

A intervenção será compreendida como uma unidade formada pela sequência didática, pela plataforma, pela mediação docente e pelos instrumentos de produção de dados. Separar artificialmente esses elementos seria inadequado, pois a plataforma só adquire o sentido previsto quando está vinculada à pergunta proposta, à previsão solicitada e à discussão posterior.

Reconhecer essa unidade não autoriza atribuir qualquer mudança observada a um único componente. As produções dos estudantes serão lidas como evidências construídas em um percurso de ensino situado, no qual tempo de aula, infraestrutura, retomadas conceituais e interação entre duplas podem interferir na forma como uma explicação aparece. A posição do professor-pesquisador exigirá registro claro das condições de aplicação e prudência na interpretação dos dados.

A definição metodológica final deverá ser discutida com o orientador antes da submissão do projeto e do início da produção de dados. Essa decisão dependerá também das condições confirmadas do campo e das exigências éticas e institucionais aplicáveis.

\section{Contexto e participantes}

O campo previsto é uma turma de Ensino Médio de escola pública na qual o pesquisador atue como professor. A possibilidade inicial é uma turma vinculada à disciplina Projeto Permanente por Afinidade, em escola de tempo integral. Esse cenário não está garantido para o próximo ano letivo, pois a lotação docente e a composição das turmas poderão mudar.

Se necessário, a sequência poderá ser aplicada em outra turma do Ensino Médio sob docência do pesquisador, inclusive em um componente de Matemática, desde que sejam preservadas as condições pedagógicas e éticas definidas para a pesquisa.

Os participantes serão estudantes iniciantes em \textit{machine learning}. Não haverá exigência de programação prévia. A sequência pressupõe leitura de tabelas e gráficos elementares, localização de pontos em eixos cartesianos, comparação de valores e mobilização de noções iniciais de média, distância, função e probabilidade quando esses conhecimentos forem retomados durante os módulos.

O número de participantes, a série, o componente curricular, a escola e o calendário de aplicação serão informados após a confirmação do campo.

\section{Etapas da intervenção}

O primeiro momento corresponderá ao estudo das referências e ao aperfeiçoamento da plataforma, dos exemplos sintéticos e dos instrumentos. Em seguida, a sequência didática será preparada de acordo com a turma e as condições de infraestrutura disponíveis.

Os estudantes trabalharão em duplas, cada uma com um computador do laboratório de informática da escola, que dispõe de 20 máquinas. O trabalho em duplas cria espaço para que os estudantes comparem hipóteses diante do mesmo experimento. Durante a aula, as previsões são individuais e depois discutidas em dupla, mas as fontes analisadas na pesquisa são apenas o diagnóstico inicial, o desafio final e os registros do professor. Isso permite distinguir o que foi discutido coletivamente do que cada estudante consegue registrar e justificar. Os oito encontros correspondem a uma duração viável para esse campo e exigirão atenção às retomadas e ao ritmo efetivo da turma.

Ao final de cada encontro, o professor-pesquisador organizará uma síntese breve que recupere os elementos discutidos pela turma.

As sínteses de encerramento não terão a função de produzir uma resposta correta única para ser copiada pelos estudantes. Elas servirão para recuperar o que foi observado, explicitar relações ainda confusas e preparar a passagem para o encontro seguinte. Caso a turma confunda uma probabilidade com garantia de ocorrência, por exemplo, a retomada poderá recuperar o papel do limiar e distinguir estimativa de certeza antes de avançar para outro modelo.

O registro docente indicará que a retomada ocorreu e quais dificuldades a motivaram, sem substituir as respostas individuais como fonte de evidência da compreensão. Essa separação ajuda a preservar o vínculo entre a mediação necessária ao ensino e a análise cuidadosa do que cada participante conseguiu elaborar.

A plataforma não registrará respostas ou dados pessoais. Os registros de pesquisa serão produzidos apenas pelos instrumentos descritos a seguir.

\section{Instrumentos de produção de dados}

Antes da sequência, os participantes realizarão diagnóstico individual breve. O instrumento buscará identificar ideias iniciais sobre dados, modelos, treinamento, previsão, classificação, agrupamento e resultado de alterações simples.

Ele não terá finalidade classificatória nem será tratado como prova de domínio prévio. Servirá para situar o ponto de partida das explicações dos estudantes.

Ao final da sequência, será aplicado desafio individual de síntese. Ele apresentará situação nova e solicitará que o estudante identifique os dados necessários, reconheça o tipo de tarefa, indique um modelo compatível e explique o que poderia ocorrer com uma mudança de dado, parâmetro ou regra.

Algumas questões serão equivalentes ou comparáveis às do diagnóstico. Isso permitirá examinar mudanças no uso dos conceitos sem reduzir a análise a uma contagem de acertos.

Depois de cada aula, o professor-pesquisador produzirá um registro breve, em tópicos, sobre dificuldades recorrentes, condições de uso, intervenções realizadas e situações que ajudem a interpretar as respostas dos estudantes. Não será solicitada uma cartilha extensa nem haverá intenção de coletar uma produção individual para cada módulo.

As perguntas inseridas no site terão função formativa imediata. Elas não serão utilizadas como fonte principal de dados, pois as respostas não permanecem salvas.

As versões preliminares do diagnóstico e do desafio final serão organizadas em instrumento próprio. Elas mobilizarão dimensões comparáveis: dados e finalidade da tarefa, previsão sobre uma alteração, interpretação de representação e possibilidade ou limite de uma decisão automatizada. O diagnóstico não exigirá nomes de modelos; o desafio final solicitará o emprego dos conceitos trabalhados quando pertinente. A comparação será qualitativa e descritiva, sem conversão automática das respostas em nota de ganho.

A comparabilidade entre os instrumentos não dependerá de repetir a mesma pergunta ou de apresentar novamente o mesmo cenário. Repetições literais podem favorecer a lembrança da atividade anterior sem revelar se a relação conceitual foi mobilizada em outra situação.

O diagnóstico poderá perguntar, por exemplo, que informações seriam necessárias para uma recomendação ou classificação simples. O desafio final apresentará um conjunto de dados diferente e solicitará que o estudante justifique a escolha de uma tarefa, de um modelo ou de uma alteração possível. A dimensão observada permanece a mesma, mas a situação exige que a resposta seja reconstruída. Quando houver item equivalente, sua formulação deverá manter o foco na explicação, não apenas no reconhecimento de uma palavra trabalhada durante a sequência.

Os instrumentos também precisarão ser compatíveis com o tempo escolar disponível e com o fato de que os participantes serão estudantes do Ensino Médio em uma atividade inserida no cotidiano da turma. Não se pretende aplicar teste extenso que interrompa a sequência ou transforme a intervenção em rotina de avaliação formal.

Perguntas abertas curtas, acompanhadas de gráfico, tabela ou situação-problema, podem oferecer indícios mais úteis que grande quantidade de itens objetivos, desde que o enunciado deixe claro o que se pede. Antes da aplicação, o professor-pesquisador verificará se cada questão está ligada a um objetivo de aprendizagem e a uma atividade efetivamente realizada.

Caso algum conceito não tenha sido abordado nas condições previstas, ele não deverá ser tratado como ausência de compreensão no desafio final. Essa precaução mantém os instrumentos vinculados ao percurso de ensino e protege a análise contra interpretações que ultrapassem as evidências disponíveis.

\section{Procedimentos de análise}

As respostas ao diagnóstico e ao desafio final serão organizadas por códigos, sem identificação nominal dos participantes. A análise será orientada por critérios construídos a partir dos objetivos de aprendizagem da sequência: diferenciar dados, características e finalidade do modelo; antecipar consequência de alteração simples em dado, parâmetro ou regra; interpretar uma representação; e justificar o resultado observado, relacionando entrada, procedimento computacional e saída com apoio no conceito ou na matemática trabalhada.

A primeira etapa consistirá na leitura integral das respostas e na organização de uma matriz por participante e dimensão de análise. Em cada dimensão, será registrada a presença, a ausência ou a forma de manifestação de uma explicação, sem apresentar essa marcação como nota individual.

Uma resposta poderá apenas indicar que ``a classificação mudou'', enquanto outra relacionará a mudança ao limiar escolhido ou à proximidade entre exemplos. A diferença entre esses modos de explicar será mais informativa que uma oposição simples entre certo e errado.

Em seguida, as respostas do diagnóstico e do desafio final serão colocadas em relação, procurando continuidades, deslocamentos e dificuldades persistentes. A análise considerará, por exemplo, se o estudante passa a distinguir uma tarefa de previsão de uma tarefa de classificação, se identifica dados compatíveis com um problema e se consegue prever uma consequência antes de observá-la na tela.

Quando houver condições para isso, frequências descritivas poderão indicar quantas produções apresentam determinada característica. Elas servirão para situar o conjunto, não para sustentar inferências estatísticas ou generalizações sobre outras turmas.

Os registros não serão apresentados integralmente na dissertação. O texto final deverá reunir sínteses das respostas, frequências descritivas quando contribuírem para a leitura do conjunto e poucos excertos anonimizados que ilustrem modos distintos de explicação. Os registros docentes contextualizarão os resultados quando uma dificuldade recorrente, uma retomada matemática ou uma condição de infraestrutura ajudar a compreender as produções escritas.

A interpretação aproximará o que o estudante escreveu das condições em que a atividade ocorreu, mas não usará registros do professor para completar ou atribuir ao participante compreensão que não esteja evidenciada em sua produção individual. Essa separação é inegociável para que a análise permaneça prudente e compatível com as fontes produzidas.

\section{Aspectos éticos}

A produção de dados somente poderá começar após a confirmação das exigências institucionais e éticas aplicáveis ao projeto. A Resolução CNS nº 510/2016 dispõe sobre pesquisas em Ciências Humanas e Sociais que utilizam dados obtidos diretamente de participantes ou informações identificáveis, enquanto a Resolução CNS nº 466/2012 estabelece diretrizes gerais para pesquisas envolvendo seres humanos \cite{brasil2016cns510,brasil2012cns466}.

Como a intervenção poderá envolver adolescentes, os procedimentos éticos deverão ser definidos conforme a orientação institucional e a tramitação exigida antes da aplicação.

As respostas serão anonimizadas por códigos, e documentos com informações de identificação, quando necessários, serão mantidos separadamente dos instrumentos de análise. Dados brutos não serão inseridos no código da plataforma nem publicados em repositório aberto. A participação será voluntária. Estudantes que não participarem da produção de dados deverão ter acesso a atividade pedagógica equivalente, sem prejuízo da experiência escolar.
